% Generated draft: requires mathematical and source review.
\documentclass[11pt]{article}
\usepackage[margin=25mm]{geometry}
\usepackage{fontspec}
\setmainfont{Latin Modern Roman}
\usepackage{amsmath,amssymb,mathtools,unicode-math}
\setmathfont{Latin Modern Math}
\usepackage{enumitem,needspace,xurl,hyperref}
\setlength{\parindent}{0pt}
\setlength{\parskip}{5pt}
\setlength{\emergencystretch}{3em}
\newcommand{\N}{\mathbb{N}}
\newcommand{\Z}{\mathbb{Z}}
\newcommand{\Q}{\mathbb{Q}}
\newcommand{\R}{\mathbb{R}}
\newcommand{\C}{\mathbb{C}}
\newcommand{\F}{\mathbb{F}}
\newcommand{\A}{\mathbb{A}}
\newcommand{\PP}{\mathbb{P}}
\newcommand{\E}{\mathbb{E}}
\newcommand{\eps}{\varepsilon}
\newcommand{\dd}{\,\mathrm{d}}
\begin{document}
\section*{UnsolvedMath ID: 10400137}
\textbf{Source problem code:} catalogue entry
\subsection{Definitions and mathematical statement}
Problem 7.22, attributed to S. Baseilhac and R. Benedetti, asks for a suitable \((2+1)\)-dimensional decorated cobordism theory supporting a quantum field theory that contains the already defined QHI. The source explicitly allows the theory to be non-purely topological. The abbreviation QHI ordinarily refers in this context to quantum hyperbolic invariants, but their input data, parameters and normalization are absent from this extract.

A decorated \((2+1)\)-cobordism framework would specify decorated surface objects and decorated \(3\)-dimensional cobordisms between them, their equivalences, and compatible composition by gluing. A field-theoretic extension would assign state objects and cobordism maps obeying the intended composition and disjoint-union rules. The relevant target category and any projective or anomaly conventions are not supplied, so they cannot be imposed by default. To contain QHI, the theory must recover the existing invariant on its original admissible closed decorated \(3\)-dimensional inputs, using the source's normalization. This requirement must allow the additional non-topological dependence mentioned in the question.

The problem particularly asks to determine the behavior of QHI under connected sums. For admissible decorated inputs \((M,\delta)\) and \((N,\eta)\), one must first define the induced connected-sum decoration and then determine the resulting invariant. No multiplicative formula, gluing normalization, or correction factor is asserted in the supplied statement. The missing definition of QHI and its decorations prevents stating such a formula in advance.
\subsection{Short English statement}
Extend the existing QHI to a decorated three-dimensional cobordism field theory, allowing dependence beyond bare topology, and especially determine how the invariants behave under connected sum. The record does not define the decorations or QHI and does not propose a connected-sum formula.
\subsection{Sources}
\begin{itemize}
\item Local UnsolvedMath registry, permanent ID \(10400137\), exact source code \texttt{catalogue entry}, Problem 7.22, attributed to S. Baseilhac and R. Benedetti. The title is truncated after ``(S.''
\item Ohtsuki, \emph{Problems on invariants of knots and 3-manifolds} (2002), Problem 7.22, PDF page 115: \url{https://msp.org/gtm/2002/04/gtm-2002-04-024s.pdf}. Bold printed headings distinguish 231 numbered items from references in surrounding exposition.
\item Extraction pdf-font-structure, accessed 2026-07-29. Source evidence NEEDS\_REVIEW presents an open problem and calls for a current primary-literature audit. Supplied research classification OPEN-TRIAGE calls the decorated-cobordism theory for QHI not obviously resolved. No separate definition or primary reference for the already defined QHI is included.
\item Dataset attribution: Ulam AI / UnsolvedMath Contributors, CC BY 4.0. Source rights NEEDS\_REVIEW: publicly accessible author/publisher copy, redistribution terms requiring review.
\end{itemize}
\end{document}
